\documentclass[11pt]{article}
\usepackage[margin=1in]{geometry}
\usepackage[T1]{fontenc}
\usepackage{lmodern}
\usepackage{amsmath,amssymb,amsthm}
\usepackage{microtype}
\usepackage{enumitem}
\usepackage{xcolor}
\usepackage[colorlinks=true,linkcolor=blue!45!black,citecolor=blue!45!black,urlcolor=blue!45!black]{hyperref}
\usepackage{booktabs}
\setlength{\parindent}{0pt}
\setlength{\parskip}{0.55em}
\setlist{itemsep=0.25em,topsep=0.35em}
\newcommand{\E}{\mathbb{E}}
\newcommand{\Prob}{\mathbb{P}}
\newcommand{\R}{\mathbb{R}}
\newtheorem{theorem}{Theorem}
\newtheorem{proposition}{Proposition}
\title{A Chance-Adjusted Normalization of PSSI\\[0.3em]
\large Motivation, aggregation consistency, and a uniqueness theorem}
\author{}
\date{}

\begin{document}
\maketitle
\vspace{-1.5em}

\section{Motivation: distinguish population composition from excess segregation}

Raw PSSI is sensitive to the population share of the racial group being studied.
A group that comprises a larger fraction of a city has more opportunities for
same-group interaction, even under random mixing. Consequently, a larger raw
score can reflect a larger population share rather than a stronger departure
from random mixing. A city that is 30\% Black should not automatically be judged
more segregated than a city that is 10\% Black merely because its raw PSSI is
higher.

For the project formulation, let $P$ be a nonnegative row-stochastic interaction
matrix, let $r_i$ be the studied group's population proportion in unit $i$, and
write
\begin{equation}
 A=P\operatorname{diag}(r),\qquad \lambda=\rho(A),
 \label{eq:raw}
\end{equation}
where $\rho$ is the spectral radius. If every unit has the same group proportion
$x$, then
\[
 A=xP,\qquad \lambda=x\rho(P)=x.
\]
Thus even perfectly uniform local composition produces a raw score of $0.10$
when $x=0.10$ and $0.30$ when $x=0.30$. This establishes a composition effect;
it does not assert that a larger group share always dominates every other
difference between two cities.

Echenique and Fryer explicitly discuss composition sensitivity and estimate a
random-allocation baseline for SSI \cite[Section VI.B, manuscript p.~25]{EF}. Here the objective
is to calibrate the chosen raw statistic against a city-specific null model
while preserving its arithmetic aggregation structure.

\section{The reference model and notation}

Fix one city and one studied group. Let $Q_0$ be a specified random-mixing model,
and let $\Lambda$ denote the raw PSSI produced by a random configuration drawn
from $Q_0$. Distinguish the observed value $\lambda$ from this random variable.
Define
\begin{equation}
 m_0=\E_0[\Lambda],\qquad 0\leq m_0<1.
 \label{eq:baseline}
\end{equation}
The desired normalization is a function $\mu:[0,1]\to\R$. Negative values are
allowed because a configuration may score below its random-mixing baseline.
All expectations and uses of $\mu$ below refer to this same fixed calibration.

One concrete null model fixes the interaction matrix $P$, the population $n_i>0$
of each unit, and the city's total group count $K$. With $N=\sum_i n_i$, assign
the $K$ group labels uniformly to the $N$ population slots. If $B_i$ is the
random group count in unit $i$, then
\[
 \Prob_0(B_1=b_1,\ldots,B_m=b_m)
 =\frac{\prod_{i=1}^{m}\binom{n_i}{b_i}}{\binom{N}{K}},
 \qquad \sum_i b_i=K.
\]
Compute $\Lambda$ using the same raw-statistic definition as for the observed
city. This preserves local population totals and the citywide group share.
The theorem below applies to any specified $Q_0$ satisfying
\eqref{eq:baseline}; this particular allocation model is a substantive choice.

Uniform composition and random allocation are different reference points.
Although $\E_0[B_i/n_i]=K/N$, in general
\[
 \E_0\!\left[\rho\!\left(P\operatorname{diag}(B_i/n_i)\right)\right]
 \ne
 \rho\!\left(P\operatorname{diag}(\E_0[B_i/n_i])\right).
\]
Accordingly, the baseline is $m_0$, not automatically the population share $K/N$.

The reference value $1$ is justified by the bound
\[
 0\leq\rho(A)\leq\|A\|_{\infty}
 =\max_i\sum_jP_{ij}r_j\leq1.
\]
It is a universal upper bound, not necessarily an attainable value for a
particular fixed city.

\section{Three axioms for the normalization}

\paragraph{Axiom 1: preservation of weighted averages.}
For every positive integer $k$, every $y_1,\ldots,y_k\in[0,1]$, and every
$w_1,\ldots,w_k\geq0$ with $\sum_iw_i=1$, require
\begin{equation}
 \boxed{\displaystyle
 \mu\!\left(\sum_{i=1}^{k}w_i y_i\right)
 =\sum_{i=1}^{k}w_i\mu(y_i).}
 \label{eq:averaging}
\end{equation}
An aggregate score should have the same normalized value whether one first
aggregates the raw scores and then normalizes, or first normalizes the scores
and then aggregates using the original weights. This preserves the
interpretation of those weights and avoids introducing an additional
preference for dispersion among component scores. The connection to SSI's
component decomposition is made precise in Section~\ref{sec:components}.

This is an axiom about averaging \emph{scores under one calibration}. It does
not claim that the spectral radius of an average of matrices equals the average
of their spectral radii, or that physically merging cities averages their scores.

\paragraph{Axiom 2: zero expectation under random mixing.}
Require
\begin{equation}
 \boxed{\E_0[\mu(\Lambda)]=0.}
 \label{eq:zero}
\end{equation}
Random configurations should have adjusted score zero on average. Applying
this requirement separately to each city removes differences in their expected
random baselines. It does not require every random configuration to receive
zero, and it does not require the full distributions of adjusted scores to be
identical across cities. The notation $\E_0[\mu(\Lambda)]$ makes explicit which
random quantity is being averaged.

\paragraph{Axiom 3: preservation of the upper reference.}
Require
\begin{equation}
 \boxed{\mu(1)=1.}
 \label{eq:upper}
\end{equation}
The universal upper reference remains one, providing a common endpoint and
fixing the scale of the adjusted index. When $1$ is unattainable under a city's
constraints, this is a calibration at a theoretical reference value. It does
not assert that the city's largest feasible adjusted score must equal one.

\section{Why weighted averages matter for SSI}
\label{sec:components}

Echenique and Fryer state:
\begin{quote}
``The SSI is the weighted average of the SSI by connected component
\ldots\ weighting each component by how many individuals it has.''
\end{quote}
\hfill\cite[footnote 9, author-posted manuscript p.~14]{EF}

Their Proposition~4 gives the decomposition into component eigenvalues.
Let $C_1,\ldots,C_k$ be the connected components of the studied group's
interaction graph, let $v_i=|C_i|$, and let $\lambda_i$ be the spectral radius
of its component matrix. In our notation, their aggregate SSI is
\begin{equation}
 S=\sum_{i=1}^{k}w_i\lambda_i,
 \qquad w_i=\frac{v_i}{\sum_{j=1}^{k}v_j}.
 \label{eq:EF}
\end{equation}
The weights count members of the studied group in each component.

\begin{proposition}[Consistency with component aggregation]
Use a single city-level normalization $\mu$ for the component scores and the
aggregate score in \eqref{eq:EF}. If Axiom~1 holds, the original component
weights also aggregate the normalized scores:
\[
 \mu(S)=\sum_{i=1}^{k}\frac{v_i}{\sum_jv_j}\,\mu(\lambda_i).
\]
\end{proposition}
\begin{proof}
The weights $w_i=v_i/\sum_jv_j$ are nonnegative and sum to one. Substitute
$S=\sum_iw_i\lambda_i$ into \eqref{eq:averaging}. Thus normalizing before or
after component aggregation produces exactly the same answer.
\end{proof}

For example, two components with weights $1/4$ and $3/4$ and raw scores $0.4$
and $0.8$ have aggregate score $0.7$. Axiom~1 requires
\[
 \mu(0.7)=\tfrac14\mu(0.4)+\tfrac34\mu(0.8).
\]
A nonlinear change of scale can violate this identity: squaring first gives
$\tfrac14(0.4)^2+\tfrac34(0.8)^2=0.52$, whereas squaring the aggregate gives
$(0.7)^2=0.49$. A positive affine calibration of the squared values leaves
this discrepancy nonzero. The adjusted aggregate would then depend on the
order of normalization and aggregation.

\paragraph{Scope of this connection.}
The same city-level $\mu$ must be used throughout. Independently calibrating
each component to its own null mean generally produces different functions
$\mu_i$ and does not preserve \eqref{eq:EF} with the original weights. In
particular, the city-centered component scores need not each have null mean
zero. Our averaging axiom is a requirement for the normalization; it is distinct
from the authors' axiom named \emph{Linearity}, which defines the underlying SSI.

Also, a single global spectral radius and the aggregate in \eqref{eq:EF}
are different constructions on a disconnected graph:
\[
 \rho\!\left(\operatorname{diag}(A_1,\ldots,A_k)\right)
 =\max_i\rho(A_i),
 \qquad
 S=\sum_i w_i\rho(A_i).
\]
The normalization theorem applies to either raw statistic in $[0,1]$. The
component-aggregation proposition applies when the chosen raw statistic
actually uses \eqref{eq:EF}; it does not turn the project's single global
spectral radius into a component average. The null simulations must use
whichever raw definition is selected.

\section{Derivation and uniqueness}

\begin{theorem}[Unique averaging-consistent, chance-adjusted normalization]
Let $\Lambda\in[0,1]$ have a specified null distribution $Q_0$ with
$m_0=\E_0[\Lambda]<1$. There is exactly one function
$\mu:[0,1]\to\R$ satisfying Axioms~1--3. It is
\begin{equation}
 \boxed{\displaystyle
 \mu(\lambda)=\frac{\lambda-m_0}{1-m_0}.}
 \label{eq:result}
\end{equation}
In particular, the axioms imply an affine transformation without separately
assuming continuity, differentiability, or monotonicity.
\end{theorem}

\begin{proof}
\textit{Necessity.}
Every $\lambda\in[0,1]$ is a weighted average of the endpoints:
\[
 \lambda=\lambda\cdot1+(1-\lambda)\cdot0.
\]
Axiom~1 therefore implies
\[
 \mu(\lambda)=\lambda\mu(1)+(1-\lambda)\mu(0).
\]
By Axiom~3, $\mu(1)=1$. Writing $b=\mu(0)$ yields
\begin{equation}
 \mu(\lambda)=\lambda+(1-\lambda)b
             =(1-b)\lambda+b.
 \label{eq:affine}
\end{equation}
Thus weighted-average consistency already forces an affine function.
Taking the null expectation and using Axiom~2 gives
\[
 0=\E_0[\mu(\Lambda)]
  =m_0+(1-m_0)b.
\]
Since $m_0<1$, this equation has the unique solution
\[
 b=-\frac{m_0}{1-m_0}.
\]
Substituting into \eqref{eq:affine},
\[
 \mu(\lambda)
 =\lambda-\frac{m_0(1-\lambda)}{1-m_0}
 =\frac{\lambda-m_0}{1-m_0}.
\]
Hence every function satisfying the axioms must have the form
\eqref{eq:result}.

\textit{Existence.}
Conversely, define $\mu$ by \eqref{eq:result}. For admissible weights,
\begin{align*}
 \mu\!\left(\sum_iw_i y_i\right)
 &=\frac{\sum_iw_i y_i-m_0}{1-m_0}\\
 &=\sum_iw_i\frac{y_i-m_0}{1-m_0}
 =\sum_iw_i\mu(y_i),
\end{align*}
where the second equality uses $\sum_iw_i=1$. Furthermore,
\[
 \E_0[\mu(\Lambda)]
 =\frac{\E_0[\Lambda]-m_0}{1-m_0}=0,
 \qquad
 \mu(1)=\frac{1-m_0}{1-m_0}=1.
\]
All three axioms hold, completing the existence and uniqueness proof.
\end{proof}

\paragraph{Interpretation and limits.}
The numerator is excess PSSI above its random-mixing mean; the denominator is
the gap from that mean to the universal upper reference. The result is
\emph{affine} in the raw score (often called a linear normalization), with
positive slope $1/(1-m_0)$. It has $\mu(m_0)=0$, preserves ordering under a
fixed calibration, and may be negative. For every city $c$ with its own null
mean $m_{0,c}<1$,
\[
 \E_{0,c}\!\left[\frac{\Lambda_c-m_{0,c}}{1-m_{0,c}}\right]=0.
\]
This is exact calibration of expected random baselines, not a proof of
population-share invariance under every segregated alternative. If $m_0=1$,
boundedness implies $\Lambda=1$ almost surely, so Axioms~2 and~3 contradict
each other; the exclusion of this case is necessary.

In computation, estimate $m_0$ by averaging PSSI over independent null
simulations and substitute that estimate in \eqref{eq:result}. The exact
theorem concerns the true expectation; a simulated mean introduces numerical
approximation error.

\begin{thebibliography}{9}
\bibitem{EF}
Federico Echenique and Roland G. Fryer, Jr. (2007).
\emph{A Measure of Segregation Based on Social Interactions}.
\emph{The Quarterly Journal of Economics}, 122(2), 441--485.
\href{https://doi.org/10.1162/qjec.122.2.441}{doi:10.1162/qjec.122.2.441}.
\href{https://eml.berkeley.edu/~fechenique/published/echenique-fryer-qje.pdf}{Author-posted manuscript}
(dated June 2006, with publication information on its cover).
The quotation is from footnote~9 on manuscript p.~14; the component formula
is Proposition~4 on manuscript p.~32. These page numbers refer to the
author-posted manuscript, not the journal pagination.
\end{thebibliography}

\end{document}
